
\subsection{Causal nets}\label{sec:causal-nets}

Causal nets are a process-mining formalism that refines Petri nets by equipping each activity with explicit input and output bindings. Where a Petri net records token multiplicity along a flow relation, a causal net records which subsets of predecessor activities jointly enable a given activity, and which subsets of successor activities are jointly produced. We translate OC causal nets into decorated cospans by the same cospan-algebra construction as for OCPNs. The only difference is that the interface sets come from bindings rather than from pre/post-set multiplicities. The reading is direct. A binding is an AND mediator, synchronising the typed pairs that fire together. The alternative bindings of an activity exclude one another, so one XOR mediator on each side of the activity chooses among them. We begin with bindings as plain sets of typed pairs. How many objects flow along each is a separate, later layer, the constraint system $\Lambda$ (Remark~\ref{rem:cnet-binding-numbers}).

\begin{definition}[Object-centric causal net]
  \label{def:causal-net}
  The object-centric (OC) causal net used here, or OCCN, is our formulation, drawn from causal
  nets~\cite{vanderaalstCausalNetsModeling2011} and their object-centric
  extension~\cite{lissObjectCentricCausalNets2025}. It is a tuple
  \[
    CN=(T,D,\mathrm{In},\mathrm{Out},\mathcal{O}),
  \]
  where \(T\) is a finite set of activities, \(\mathcal{O}\) is a finite set of object types, and
  \[
    D\subseteq (T\cup\{\bot\})\times\mathcal{O}\times (T\cup\{\top\})
  \]
  is a typed dependency multigraph, where each arc $(s,\omega,t)\in D$ records that $s$ can produce a token of type $\omega\in\mathcal{O}$ consumed by $t$. The symbols $\bot,\top\notin T$ are the open ends of Definition~\ref{def:contexts-general}, where objects enter and leave the net. The maps
  \begin{gather*}
    \mathrm{In}:T\to \mathcal{P}(\mathcal{P}((T\cup\{\bot\})\times\mathcal{O})),\\
    \mathrm{Out}:T\to \mathcal{P}(\mathcal{P}((T\cup\{\top\})\times\mathcal{O}))
  \end{gather*}
  assign to each activity $t$ a set of input bindings and a set of output bindings. Each input binding $X\in \mathrm{In}(t)$ is a finite set of $(s,\omega)$-pairs (typed predecessors jointly enabling $t$), and each output binding $Y\in \mathrm{Out}(t)$ is a finite set of $(t',\omega)$-pairs (typed successors jointly produced when $t$ fires), all consistent with~$D$.
\end{definition}

This construction instantiates the general framework of Section~\ref{sec:general-framework}, and
the mapping is the most direct of the four. An activity is a plain node, and each binding is an
AND mediator. An activity's several input (or output) bindings are exclusive alternatives, so one
XOR mediator on each side of the activity has the AND mediators of its bindings as branches. A
singleton binding is a one-in, one-out AND that passes flow straight through. A pair $(\bot,\omega)$
or $(\top,\omega)$ in a binding gives its AND mediator an edge from a mediator with no in-edges or
to one with no out-edges, so the walk ends there at the open end $\bot$ or $\top$
(Definition~\ref{def:contexts-general}). This determines the instance map.
\medskip
\begin{center}
  \small\begin{tabularx}{\linewidth}{@{}lX@{}}
    \hline
    General framework & Causal net \\
    \hline
    AND/OR graph $(V,E)$ & the activities $T$, with one XOR mediator on each side of an activity and one AND mediator per binding; typed wires read from $D$ \\
    Activities $\mathcal{A}$ & the activities $T$ \\
    AND mediators & one per binding in $\mathrm{In}(t),\mathrm{Out}(t)$, synchronising the typed pairs that fire together (a singleton binding is a pass-through) \\
    XOR mediators & one on each side of each activity $t$, choosing among the bindings of $\mathrm{In}(t)$ or of $\mathrm{Out}(t)$ \\
    \hline
  \end{tabularx}
\end{center}
\medskip
Running the construction of Definition~\ref{def:contexts-general} through this mapping is
immediate, since the bindings are already the contexts. The walk from an activity $t$ steps
through one XOR and one AND mediator on each side and reaches its neighbours at once, so its
backward and forward families are the input and output bindings themselves, $\mathrm{In}(t)$ and
$\mathrm{Out}(t)$. Every pair of an input binding with an output binding is a context,
\[
  \mathrm{Contexts}_{CN}(t):=\mathrm{In}(t)\times \mathrm{Out}(t),
\]
and each context $c=(P,S)$ yields one generator with typed wires
$L_c=\{(s,\{\omega\},t):(s,\omega)\in P\}$ and $R_c=\{(t,\{\omega\},t'):(t',\omega)\in S\}$.
Counts are carried separately by the constraint system $\Lambda$
(Remark~\ref{rem:cnet-binding-numbers}), so a split of one type across two output legs, as at the
review $r$ of Example~\ref{ex:running-cnet-cospan} below, is a context like any other. Generators
sharing a typed wire compose along it in $\mathbf{F}(\Sigma)$ (Definition~\ref{def:free-category}).
Start and end generators attached at $\bot$ and $\top$ belong to $\Sigma_\gamma$, a choice of the
analysis (Section~\ref{subsec:analysis}).

\begin{theorem}[Canonical cospan-algebra presentation of OC causal nets]
  \label{thm:causal-to-cospan-signature}
  Let $CN = (T, D, \mathrm{In}, \mathrm{Out}, \mathcal{O})$ be an OC causal net. Then $CN$ canonically determines a cospan-algebra signature $\Sigma$, and hence the free symmetric monoidal category $\mathbf{F}(\Sigma)$.
\end{theorem}

\begin{proof}
  The schema of Section~\ref{sec:notation-by-notation} applies at the instance map above,
  finiteness of $T$ and of all binding sets giving the conditions of
  Definition~\ref{def:lm-graph}, with each edge typed by the arc types in~$D$. The delta is that
  the walk is immediate. It steps through one XOR and one AND mediator on each side of an activity and
  reaches its neighbours at once, returning $\mathrm{In}(t)$ and $\mathrm{Out}(t)$. Every pair
  of them is a context, and
  $\mathrm{Contexts}_{CN}(t)=\mathrm{In}(t)\times \mathrm{Out}(t)$ is finite because $T$ and
  the binding sets are finite.
\end{proof}

\begin{figure*}[tp]
  \centering
  \begin{tikzpicture}[scale=1.0, inline,
    cnact/.style={rectangle, rounded corners=2pt, draw=black!55, fill=black!5,
                  minimum width=7mm, minimum height=12mm, font=\small},
    bnd/.style={rectangle, rounded corners=2pt, draw=black!55, fill=black!12,
                minimum width=7mm, minimum height=12mm, font=\small},
    omk/.style={circle, draw=#1, fill=#1, minimum size=4.6pt, inner sep=0pt},
    smk/.style={rectangle, sharp corners, draw=#1, fill=#1, minimum size=5.4pt, inner sep=0pt},
    card/.style={font=\small, inner sep=0.5pt},
    klab/.style={font=\tiny, black!70, inner sep=0.5pt},
    andconn/.style={black!75, line width=0.7pt},
    keyconn/.style={black!75, line width=0.7pt}]   %
  \def\dy{1.8mm}
  \def\rgap{4pt}
  \node[bnd]   (g1) at (0,0)      {$\bot$};
  \node[cnact] (s)  at (2.5,0)    {$s$};
  \node[cnact] (r)  at (5.0,0)    {$r$};
  \node[cnact] (i)  at (8.2,1.30) {$i$};
  \node[cnact] (n)  at (8.2,-1.30){$n$};
  \node[bnd]   (g2) at (11.4,0)   {$\top$};

  \draw[blueflow] ([yshift=\dy]g1.east) --
      node[smk=bluecol, pos=0.25] (a1){} node[smk=bluecol, pos=0.75] (a2){} ([yshift=\dy]s.west);
  \draw[redflow]  ([yshift=-\dy]g1.east) --
      node[smk=redcol, pos=0.25] (a3){} node[smk=redcol, pos=0.75] (a4){} ([yshift=-\dy]s.west);
  \draw[andconn] (a1)--(a3);   \draw[andconn] (a2)--(a4);
  \draw[blueflow] ([yshift=\dy]s.east) --
      node[smk=bluecol, pos=0.25] (b1){}
      node[smk=bluecol, pos=0.75, label={[card,bluecol,label distance=2pt]above:$[1,5]$}] (b2){}
      ([yshift=\dy]r.west);
  \draw[redflow]  ([yshift=-\dy]s.east) --
      node[smk=redcol, pos=0.25] (b3){}
      node[smk=redcol, pos=0.75, label={[card,redcol,label distance=2pt]below:$[2,4]$}] (b4){}
      ([yshift=-\dy]r.west);
  \draw[andconn] (b1)--(b3);   \draw[andconn] (b2)--(b4);

  \draw[blueflow] ([yshift=0.54cm]r.east) --
      node[omk=bluecol, pos=0.25] (c1){} node[omk=bluecol, pos=0.75] (d1){} ([yshift=\dy]i.west);   %
  \draw[redflow]  ([yshift=0.18cm]r.east) --
      node[omk=redcol, pos=0.25] (c2){} node[omk=redcol, pos=0.75] (d2){} ([yshift=-\dy]i.west);    %
  \draw[blueflow] ([yshift=-0.18cm]r.east) --
      node[smk=bluecol, pos=0.25] (c3){} node[smk=bluecol, pos=0.75] (e1){} ([yshift=\dy]n.west);   %
  \draw[redflow]  ([yshift=-0.54cm]r.east) --
      node[smk=redcol, pos=0.25] (c4){} node[smk=redcol, pos=0.75] (e2){} ([yshift=-\dy]n.west);    %
  \draw[keyconn] (c1)--(c2)--(c3)--(c4);            %
  \node[klab, anchor=south west] at ([shift={(1pt,\rgap)}]c1) {$\rho_{o}$};  \node[klab, anchor=south west] at ([shift={(1pt,\rgap)}]c3) {$\rho_{o}$};
  \node[klab, anchor=north west] at ([shift={(1pt,-\rgap)}]c2) {$\rho_{it}$}; \node[klab, anchor=north west] at ([shift={(1pt,-\rgap)}]c4) {$\rho_{it}$};
  \draw[andconn] (d1)--(d2);   %
  \draw[andconn] (e1)--(e2);   %

  \draw[blueflow] ([yshift=\dy]i.east) --
      node[omk=bluecol, pos=0.25] (f1){} node[omk=bluecol, pos=0.75] (h1){} ([yshift=0.54cm]g2.west);
  \draw[redflow]  ([yshift=-\dy]i.east) --
      node[omk=redcol, pos=0.25] (f2){} node[omk=redcol, pos=0.75] (h3){} ([yshift=0.18cm]g2.west);
  \draw[andconn] (f1)--(f2);   %
  \draw[blueflow] ([yshift=\dy]n.east) --
      node[smk=bluecol, pos=0.25] (k1){} node[smk=bluecol, pos=0.75] (h5){} ([yshift=-0.18cm]g2.west);
  \draw[redflow]  ([yshift=-\dy]n.east) --
      node[smk=redcol, pos=0.25] (k2){} node[smk=redcol, pos=0.75] (h7){} ([yshift=-0.54cm]g2.west);
  \draw[andconn] (k1)--(k2);   %
  \draw[andconn] (h1)--(h3)--(h5)--(h7);   %
\end{tikzpicture}
\unskip
  \caption{An object-centric causal net over two object types, order (blue) and item (red).
    Every wire carries objects of one type, and every activity binds both types on each
    side. The solid connector joins the markers of a single binding, whose objects are
    consumed and produced together. On a wire a $\square$ marker is a batch of objects of
    that type and a $\circ$ marker is exactly one. The two bounds on the wires from the
    send $s$ to the review $r$ are declared per leg and independently of each other,
    $[1,5]$ orders and $[2,4]$ items. The review $r$ splits each type by a key of its own,
    $\rho_{o}$ over orders and $\rho_{it}$ over items. Each order and each item, taken singly,
    makes an exclusive choice between the legs that share a key, so each object leaves on one of
    them and the two legs hold the batch between them. The outcome $i$ takes one order and one
    item, and the outcome $n$ takes the remainder of each type. The boundary nodes $\bot$ and
    $\top$ are not activities of the net. They are the open ends where objects enter and leave
    (Definition~\ref{def:causal-net}).}
  \label{fig:re-occn}
\end{figure*}

\begin{figure*}[tp]
  \centering
  \begin{tikzpicture}[inline,
    bspd/.style={spider, fill=bluecol},
    rspd/.style={spider, fill=redcol},
    mbox/.style={draw=black, rounded corners=4pt, fill=white, minimum size=0.75cm, font=\small},
    plabel/.style={font=\tiny},
    cardb/.style={font=\footnotesize, bluecol, inner sep=0.5pt},
    cardr/.style={font=\footnotesize, redcol, inner sep=0.5pt}]

  \node[font=\footnotesize] at (2.7,1.9) {decorated cospan};
  \node[font=\footnotesize] at (9.0,1.9) {string diagram};

  \node[font=\small] at (-2.3,0) {$g_r$};
  \node[cospanblue, minimum height=1.0cm] (rLb) at (0,0) {};
  \node[bspd] (rLo)  at (0,0.22) {};
  \node[rspd] (rLit) at (0,-0.22) {};
  \node at (0.9,0) {$\rightarrow$};
  \node[cospangrey, minimum height=2.0cm, minimum width=2.6cm] (rA) at (2.7,0) {};
  \node[mbox, minimum height=1.0cm] (rf) at (2.7,0) {$r$};
  \node[bspd] (rao) at (1.65,0.22) {};
  \node[rspd] (rat) at (1.65,-0.22) {};
  \node[bspd] (roi) at (3.75,0.72) {};
  \node[bspd] (ron) at (3.75,0.30) {};
  \node[rspd] (rii) at (3.75,-0.30) {};
  \node[rspd] (rin) at (3.75,-0.72) {};
  \draw[bluewire] (rao) -- (rf.170);
  \draw[redwire]  (rat) -- (rf.190);
  \draw[bluewire] (rf.40)  -- (roi);
  \draw[bluewire] (rf.15)  -- (ron);
  \draw[redwire]  (rf.-15) -- (rii);
  \draw[redwire]  (rf.-40) -- (rin);
  \node at (4.5,0) {$\leftarrow$};
  \node[cospanblue, minimum height=1.9cm] (rRb) at (5.4,0) {};
  \node[bspd] (rRoi) at (5.4,0.72) {};
  \node[bspd] (rRon) at (5.4,0.30) {};
  \node[rspd] (rRii) at (5.4,-0.30) {};
  \node[rspd] (rRin) at (5.4,-0.72) {};
  \node[plabel, anchor=east] at ([shift={(-3pt,0.22cm)}]rLb.west) {$(s,\mathrm{o},r)$};
  \node[plabel, anchor=east] at ([shift={(-3pt,-0.22cm)}]rLb.west) {$(s,\mathrm{it},r)$};
  \node[plabel, anchor=west] at ([shift={(3pt,0.72cm)}]rRb.east) {$(r,\mathrm{o},i)$};
  \node[plabel, anchor=west] at ([shift={(3pt,0.30cm)}]rRb.east) {$(r,\mathrm{o},n)$};
  \node[plabel, anchor=west] at ([shift={(3pt,-0.30cm)}]rRb.east) {$(r,\mathrm{it},i)$};
  \node[plabel, anchor=west] at ([shift={(3pt,-0.72cm)}]rRb.east) {$(r,\mathrm{it},n)$};
  \node at (7.4,0) {$\mapsto$};
  \node[morphismblue, minimum height=1.6cm] (rsd) at (9.0,0) {$r$};
  \draw[bluewire] ($(rsd.south west)!0.70!(rsd.north west)$) -- ++(-0.55,0);
  \draw[redwire]  ($(rsd.south west)!0.30!(rsd.north west)$) -- ++(-0.55,0);
  \coordinate (p1) at ($($(rsd.south west)!0.70!(rsd.north west)$)-(0.275,0)$);
  \coordinate (p2) at ($($(rsd.south west)!0.30!(rsd.north west)$)-(0.275,0)$);
  \draw[gray!60, dashed] (p1) -- (p2);
  \foreach \m in {p1,p2}
    \node[circle, draw=gray!65, fill=gray!25, minimum size=3.2pt, inner sep=0] at (\m) {};
  \draw[bluewire] ($(rsd.south east)!0.82!(rsd.north east)$) -- ++(0.55,0);
  \draw[bluewire] ($(rsd.south east)!0.61!(rsd.north east)$) -- ++(0.55,0);
  \draw[redwire]  ($(rsd.south east)!0.39!(rsd.north east)$) -- ++(0.55,0);
  \draw[redwire]  ($(rsd.south east)!0.18!(rsd.north east)$) -- ++(0.55,0);
  \coordinate (m1) at ($($(rsd.south east)!0.82!(rsd.north east)$)+(0.275,0)$);
  \coordinate (m2) at ($($(rsd.south east)!0.61!(rsd.north east)$)+(0.275,0)$);
  \coordinate (m3) at ($($(rsd.south east)!0.39!(rsd.north east)$)+(0.275,0)$);
  \coordinate (m4) at ($($(rsd.south east)!0.18!(rsd.north east)$)+(0.275,0)$);
  \draw[gray!60, dashed] (m1) -- (m4);
  \foreach \m in {m1,m2,m3,m4}
    \node[circle, draw=gray!65, fill=gray!25, minimum size=3.2pt, inner sep=0] at (\m) {};
  \node[gray!75, font=\small] at ($(rsd.north)+(0,0.30)$) {$\Lambda_r$};
  \node[draw=black!45, rounded corners=2pt, fill=black!3, align=left, font=\scriptsize,
        inner sep=4pt, anchor=west] (Lam) at (-0.4,-2.05)
    {\textbf{constraint system} $\Lambda_r$\\[2pt]
     \ \ \emph{independent}\ \ $1\le n_{(s,\mathrm{o},r)}\le 5$,\ \ $2\le n_{(s,\mathrm{it},r)}\le 4$\\[1pt]
     \ \ \emph{coupling}\ \ $n_{(r,\mathrm{o},i)}+n_{(r,\mathrm{o},n)}=n_{(s,\mathrm{o},r)}$,\ \ \
     $n_{(r,\mathrm{it},i)}+n_{(r,\mathrm{it},n)}=n_{(s,\mathrm{it},r)}$};

  \begin{scope}[yshift=-3.8cm]
    \node[font=\small] at (-2.3,0) {$g_i$};
    \node[cospanblue, minimum height=0.9cm] (iL) at (0,0) {};
    \node[bspd] (iLo)  at (0,0.20) {};
    \node[rspd] (iLit) at (0,-0.20) {};
    \node at (0.9,0) {$\rightarrow$};
    \node[cospangrey, minimum height=1.0cm, minimum width=2.6cm] (iA) at (2.7,0) {};
    \node[mbox, minimum height=0.8cm] (iff) at (2.7,0) {$i$};
    \node[bspd] (iao) at (1.65,0.20) {};
    \node[rspd] (iat) at (1.65,-0.20) {};
    \node[bspd] (ibo) at (3.75,0.20) {};
    \node[rspd] (ibt) at (3.75,-0.20) {};
    \draw[bluewire] (iao) -- (iff.160);
    \draw[redwire]  (iat) -- (iff.200);
    \draw[bluewire] (iff.20)  -- (ibo);
    \draw[redwire]  (iff.-20) -- (ibt);
    \node at (4.5,0) {$\leftarrow$};
    \node[cospanblue, minimum height=0.9cm] (iR) at (5.4,0) {};
    \node[bspd] (iRo)  at (5.4,0.20) {};
    \node[rspd] (iRit) at (5.4,-0.20) {};
    \node[plabel, anchor=east] at ([shift={(-3pt,0.20cm)}]iL.west) {$(r,\mathrm{o},i)$};
    \node[plabel, anchor=east] at ([shift={(-3pt,-0.20cm)}]iL.west) {$(r,\mathrm{it},i)$};
    \node[plabel, anchor=west] at ([shift={(3pt,0.20cm)}]iR.east) {$(i,\mathrm{o},\top)$};
    \node[plabel, anchor=west] at ([shift={(3pt,-0.20cm)}]iR.east) {$(i,\mathrm{it},\top)$};
    \node at (7.4,0) {$\mapsto$};
    \node[morphismblue, minimum height=0.95cm] (isd) at (9.0,0) {$i$};
    \draw[bluewire] ($(isd.south west)!0.70!(isd.north west)$) -- ++(-0.55,0);
    \draw[redwire]  ($(isd.south west)!0.30!(isd.north west)$) -- ++(-0.55,0);
    \draw[bluewire] ($(isd.south east)!0.70!(isd.north east)$) -- ++(0.55,0);
    \draw[redwire]  ($(isd.south east)!0.30!(isd.north east)$) -- ++(0.55,0);
    \coordinate (isdp1) at ($($(isd.south west)!0.70!(isd.north west)$)-(0.275,0)$);
    \coordinate (isdp2) at ($($(isd.south west)!0.30!(isd.north west)$)-(0.275,0)$);
    \coordinate (isdm1) at ($($(isd.south east)!0.70!(isd.north east)$)+(0.275,0)$);
    \coordinate (isdm2) at ($($(isd.south east)!0.30!(isd.north east)$)+(0.275,0)$);
    \draw[gray!60, dashed] (isdp1) -- (isdp2);
    \draw[gray!60, dashed] (isdm1) -- (isdm2);
    \foreach \m in {isdp1,isdp2,isdm1,isdm2}
      \node[circle, draw=gray!65, fill=gray!25, minimum size=3.2pt, inner sep=0] at (\m) {};
    \node[gray!75, font=\small] at ($(isd.north)+(0,0.30)$) {$\Lambda_i$};
  \end{scope}

  \begin{scope}[yshift=-7.0cm]
    \node[font=\small] at (-2.3,0) {$g_n$};
    \node[cospanblue, minimum height=0.9cm] (nL) at (0,0) {};
    \node[bspd] (nLo)  at (0,0.20) {};
    \node[rspd] (nLit) at (0,-0.20) {};
    \node at (0.9,0) {$\rightarrow$};
    \node[cospangrey, minimum height=1.0cm, minimum width=2.6cm] (nA) at (2.7,0) {};
    \node[mbox, minimum height=0.8cm] (nf) at (2.7,0) {$n$};
    \node[bspd] (nao) at (1.65,0.20) {};
    \node[rspd] (nat) at (1.65,-0.20) {};
    \node[bspd] (nbo) at (3.75,0.20) {};
    \node[rspd] (nbt) at (3.75,-0.20) {};
    \draw[bluewire] (nao) -- (nf.160);
    \draw[redwire]  (nat) -- (nf.200);
    \draw[bluewire] (nf.20)  -- (nbo);
    \draw[redwire]  (nf.-20) -- (nbt);
    \node at (4.5,0) {$\leftarrow$};
    \node[cospanblue, minimum height=0.9cm] (nR) at (5.4,0) {};
    \node[bspd] (nRo)  at (5.4,0.20) {};
    \node[rspd] (nRit) at (5.4,-0.20) {};
    \node[plabel, anchor=east] at ([shift={(-3pt,0.20cm)}]nL.west) {$(r,\mathrm{o},n)$};
    \node[plabel, anchor=east] at ([shift={(-3pt,-0.20cm)}]nL.west) {$(r,\mathrm{it},n)$};
    \node[plabel, anchor=west] at ([shift={(3pt,0.20cm)}]nR.east) {$(n,\mathrm{o},\top)$};
    \node[plabel, anchor=west] at ([shift={(3pt,-0.20cm)}]nR.east) {$(n,\mathrm{it},\top)$};
    \node at (7.4,0) {$\mapsto$};
    \node[morphismblue, minimum height=0.95cm] (nsd) at (9.0,0) {$n$};
    \draw[bluewire] ($(nsd.south west)!0.70!(nsd.north west)$) -- ++(-0.55,0);
    \draw[redwire]  ($(nsd.south west)!0.30!(nsd.north west)$) -- ++(-0.55,0);
    \draw[bluewire] ($(nsd.south east)!0.70!(nsd.north east)$) -- ++(0.55,0);
    \draw[redwire]  ($(nsd.south east)!0.30!(nsd.north east)$) -- ++(0.55,0);
    \coordinate (nsdp1) at ($($(nsd.south west)!0.70!(nsd.north west)$)-(0.275,0)$);
    \coordinate (nsdp2) at ($($(nsd.south west)!0.30!(nsd.north west)$)-(0.275,0)$);
    \coordinate (nsdm1) at ($($(nsd.south east)!0.70!(nsd.north east)$)+(0.275,0)$);
    \coordinate (nsdm2) at ($($(nsd.south east)!0.30!(nsd.north east)$)+(0.275,0)$);
    \draw[gray!60, dashed] (nsdp1) -- (nsdp2);
    \draw[gray!60, dashed] (nsdm1) -- (nsdm2);
    \foreach \m in {nsdp1,nsdp2,nsdm1,nsdm2}
      \node[circle, draw=gray!65, fill=gray!25, minimum size=3.2pt, inner sep=0] at (\m) {};
    \node[gray!75, font=\small] at ($(nsd.north)+(0,0.30)$) {$\Lambda_n$};
  \end{scope}
  \node[draw=black!45, rounded corners=2pt, fill=black!3, align=left, font=\scriptsize,
        inner sep=4pt, anchor=west] (Lami) at (-0.4,-5.3)
    {\textbf{constraint system} $\Lambda_i$\\[2pt]
     \ \ \emph{independent}\ \ $n_{(r,\mathrm{o},i)}=1$,\ \ $n_{(r,\mathrm{it},i)}=1$\\[1pt]
     \ \ \emph{coupling}\ \ $n_{(i,\mathrm{o},\top)}=n_{(r,\mathrm{o},i)}$,\ \ \ $n_{(i,\mathrm{it},\top)}=n_{(r,\mathrm{it},i)}$};
  \node[draw=black!45, rounded corners=2pt, fill=black!3, align=left, font=\scriptsize,
        inner sep=4pt, anchor=west] (Lamn) at (-0.4,-8.5)
    {\textbf{constraint system} $\Lambda_n$\\[2pt]
     \ \ \emph{independent}\ \ $n_{(r,\mathrm{o},n)}\ge 1$,\ \ $n_{(r,\mathrm{it},n)}\ge 1$\\[1pt]
     \ \ \emph{coupling}\ \ $n_{(n,\mathrm{o},\top)}=n_{(r,\mathrm{o},n)}$,\ \ \ $n_{(n,\mathrm{it},\top)}=n_{(r,\mathrm{it},n)}$};
\end{tikzpicture}
\unskip
  \caption{The generators $g_r$, $g_i$ and $g_n$ of Fig.~\ref{fig:re-occn}, each shown as a decorated cospan (left)
    and the string diagram it denotes (right). Order wires are blue and item wires red, and the
    labels beside each boundary name its typed ports in the order drawn. The outcome
    generators are $g_i$, one investigated order and one item, and $g_n$, the remainder of
    each type. Each generator has a constraint system of its own, set out below its row and
    flagged on the string diagram by the $\Lambda$ marker spanning both boundaries. A
    system's independent line records the bounds that generator declares on its own legs,
    and its coupling line relates its two boundaries. Composition identifies the variables of
    shared wires and conjoins the systems, so three systems that each hold on their own need not
    hold together.}
  \label{fig:re-occn-cospans}
\end{figure*}

\begin{exmp}[A small OCCN and its generator cospans]
  \label{ex:running-cnet-cospan}
  Consider the OC causal net of Fig.~\ref{fig:re-occn}, over two object types,
  $\mathrm{order}$ and $\mathrm{item}$, each key-distributed at the review $r$. It has no places. Write $(x,\mathrm{o},y)$ for $(x,\{\mathrm{o}\},y)$. Each
  activity carries its bindings, which are already its contexts:
  \begin{gather*}
    \mathrm{Out}(s)=\{\{(r,\mathrm{o}),(r,\mathrm{it})\}\},\qquad\twocolbreak
    \mathrm{In}(r)=\{\{(s,\mathrm{o}),(s,\mathrm{it})\}\},
  \end{gather*}
  \begin{gather*}
    \mathrm{Out}(r)=\{\{(i,\mathrm{o}),(i,\mathrm{it}),\twocolbreak(n,\mathrm{o}),(n,\mathrm{it})\}\},\\
    \mathrm{In}(i)=\mathrm{In}(n)=\{\{(r,\mathrm{o}),(r,\mathrm{it})\}\}.
  \end{gather*}
  Activity $r$ has a single output binding $\{(i,\mathrm{o}),(i,\mathrm{it}),(n,\mathrm{o}),(n,\mathrm{it})\}$. It
  emits both legs on every firing and is a key-distribution split, not an
  exclusive choice. Reading a generator off each binding
  (Fig.~\ref{fig:re-occn-cospans}) gives, for $r$,
  \begin{gather*}
    g_r = \Bigl(\{(s,\mathrm{o},r),(s,\mathrm{it},r)\}
      \xrightarrow{\iota} A \xleftarrow{\iota}\twocolbreak 
      \{(r,\mathrm{o},i),(r,\mathrm{it},i),(r,\mathrm{o},n),(r,\mathrm{it},n)\},\; d_r\Bigr),
  \end{gather*}
  whose decoration carries the constraint system $\Lambda$
  (Remark~\ref{rem:cnet-binding-numbers}). Read off $g_r$'s boundary, the markings say what each
  wire carries in one firing of $r$. The input wire $(s,\mathrm{o},r)$ takes a batch of one to
  five orders ($1\le n_{(s,\mathrm{o},r)}\le 5$), and the investigated leg $(r,\mathrm{o},i)$
  holds exactly one ($n_{(r,\mathrm{o},i)}=1$). The two output legs hold the whole batch between
  them, tied by a within-type key $\rho_{o}$ that partitions it,
  $n_{(r,\mathrm{o},i)}+n_{(r,\mathrm{o},n)}=n_{(s,\mathrm{o},r)}$. Hence one firing of $r$
  receives a batch and splits it, one order to $i$ and the rest to $n$, and a batch of more than five
  forces $r$ to fire again. The item type is split in the same way by its own within-type key $\rho_{it}$, one item
  investigated on $i$ and the rest on $n$, so both types follow the identical key-distribution
  pattern under bounds of their own.
\end{exmp}

\begin{remark}[Binding numbers and keys as the constraint system]
  \label{rem:cnet-binding-numbers}
  Definition~\ref{def:causal-net} records which typed pairs bind, but not
  \emph{how many} objects flow along each. The full OC causal-net
  formalism~\cite{lissObjectCentricCausalNets2025} additionally equips every
  binding pair with a cardinality $[c_{\min},c_{\max}]$ and may correlate
  pairs through a shared object key $\rho$. Both enter the constraint system of
  Definition~\ref{def:multiplicity-decoration}. The binding cardinality is a
  per-leg interval, an inequality of $\rho(t)$ on a variable leg, and the shared key is an equality of $\rho(t)$ at its activity $t$ (Definition~\ref{def:lm-graph})
  among the legs of one generator, a partition of the input wire across the output wires it feeds.
  The routing
  structure, which generators exist, is untouched. The numbers decorate those
  generators, and an unannotated binding defaults to $n_p=1$, the unannotated sense of the
  generators of Example~\ref{ex:running-cnet-cospan}.

  On the small example above this attaches as follows. The send activity $s$ ships
  a batch of orders and a batch of items to $r$ in one firing, so each wire is bounded on its own,
  $1\le n_{(s,\mathrm{o},r)}\le 5$ for orders and $2\le n_{(s,\mathrm{it},r)}\le 4$ for items. The receive activity $r$
  then distributes each type between the two outcomes by its own within-type key. The order key
  $\rho_{o}$ partitions the orders across the investigated and not-investigated legs,
  $n_{(r,\mathrm{o},i)}+n_{(r,\mathrm{o},n)}-n_{(s,\mathrm{o},r)}=0$, with exactly one investigated
  per firing, $n_{(r,\mathrm{o},i)}=1$. The item key $\rho_{it}$ partitions the items the
  same way, $n_{(r,\mathrm{it},i)}+n_{(r,\mathrm{it},n)}-n_{(s,\mathrm{it},r)}=0$ with
  $n_{(r,\mathrm{it},i)}=1$, so the two key constraints sit side by side in $\Lambda$ and make the
  distribution explicit. Hence the number of investigations equals the number of $r$ firings, and a
  batch above either bound forces $r$ to fire again.

  Generators sharing a wire identify its leg variable and conjoin their
  constraints, and a closed composite exists if and only if the glued constraint system is
  solvable over $\mathbb{N}$. The symbolic signature carrying
  $\Lambda$ is bound-independent and is the form compared across notations.
\end{remark}

\begin{remark}[Causal net states]
  The state of a causal net mid-execution is the causal-net version of a Petri-net
  marking as a cut (Remark~\ref{rem:pn-marking-cut}). A cut through the string diagram
  carries the active wire ports, the bindings in progress rather than tokens in places,
  and sweeping the cut from boundary to boundary reads off the state evolution with no
  separate algebraic state machinery. On the OCCN example, after the $r$ generator fires
  the cut carries the four ports leading to $i$ and~$n$ simultaneously, an order and an item wire to each, and firing
  $i$ or $n$ first advances that outcome's two ports to $\top$, two binding sequences of one
  diagram.
\end{remark}
